\documentclass[runningheads]{llncs}
%
\usepackage[T1]{fontenc}
% T1 fonts will be used to generate the final print and online PDFs,
% so please use T1 fonts in your manuscript whenever possible.
% Other font encodings may result in incorrect characters.
\usepackage{cite}
\usepackage{amsmath,amssymb,amsfonts}

\usepackage{textcomp}
\usepackage[utf8]{inputenc}
\usepackage{lmodern}
\usepackage{hyperref}
\usepackage[english]{babel}
\usepackage{todonotes}
\usepackage{algorithm}
\usepackage{algpseudocode}
\usepackage{mathtools}
\usepackage{mathdots}
\usepackage{tikz-cd} 
\usepackage{wrapfig}
\usepackage{float}
\usepackage{quiver}
\usepackage{adjustbox}
\usepackage{stackengine,scalerel}
\newcommand\eye{\ensurestackMath{\stackinset{c}{}{c}{-.33pt}%
        {\bullet}{\bigcirc}}}
\newcommand\smallereye{\scalerel*{\eye}{x}}
\makeatletter
\newcommand{\removelatexerror}{\let\@latex@error\@gobble}
\makeatother

\newcounter{cases}
\newcounter{subcases}[cases]
\newenvironment{mycase}
{
    \setcounter{cases}{0}
    \setcounter{subcases}{0}
    \newcommand{\case}
    {
        \par\indent\stepcounter{cases}\textbf{Case \thecases.}
    }
    \newcommand{\subcase}
    {
        \par\indent\stepcounter{subcases}\textbf{ \>\>\>\> Subcase (\thesubcases)}
    }
}
{
    \par
}
\renewcommand*\thecases{\arabic{cases}}
\renewcommand*\thesubcases{\alph{subcases}}


\DeclareMathOperator*{\argmin}{argmin}
\DeclareMathOperator*{\pre}{^\bullet}
%\usepackage{amsthm}
\usepackage{tikz}
\usetikzlibrary{cd}
\usetikzlibrary{petri,positioning}

\spnewtheorem*{lemma*}{Lemma}{\normalshape\bfseries}{\itshape}
\spnewtheorem*{theorem*}{Theorem}{\normalshape\bfseries}{\itshape}

\usepackage{graphicx}
% Used for displaying a sample figure. If possible, figure files should
% be included in EPS format.
%
% If you use the hyperref package, please uncomment the following two lines
% to display URLs in blue roman font according to Springer's eBook style:
\usepackage{color}
\renewcommand\UrlFont{\color{blue}\rmfamily}
%
\begin{document}
%
\title{Timed Alignments with Mixed Moves
        \thanks{Neha Rino was funded by the International Master's Scholarships Program IDEX of Université Paris-Saclay.}
}%
%\titlerunning{Abbreviated paper title}
% If the paper title is too long for the running head, you can set
% an abbreviated paper title here
%
\author{Neha Rino\inst{1}\orcidID{0009-0001-3273-1954} \and
Thomas Chatain\inst{2}\orcidID{0000-0002-1470-5074}}
% First names are abbreviated in the running head.
% If there are more than two authors, 'et al.' is used.
%
\institute{Department of Computer Science, University of Warwick, Coventry, UK
  \email{Neha.Rino@warwick.ac.uk}
\and
LMF, ENS Paris-Saclay, CNRS, Université Paris-Saclay, Inria,
Gif-sur-Yvette, France
\email{thomas.chatain@ens-paris-saclay.fr} }

%
\maketitle              % typeset the header of the contribution
%% \todo{check Neha's affiliation}
%
\begin{abstract}
We study conformance checking for timed models, that is, process models that consider both the sequence of events that occur, as well as the timestamps at which each event is recorded. 
Time-aware process mining is a growing subfield of research, and as tools that seek to discover timing-related properties in processes develop, so does the need for conformance-checking techniques that can tackle time 
constraints and provide insightful quality measures for time-aware process models. 
One of the most useful conformance artefacts is the alignment, that is, finding the minimal changes necessary to correct a new observation to conform to a process model.
In this paper, we extend the notion of timed distance from a previous work where an edit on an event’s timestamp came in two types, depending on whether or not it would propagate to its successors. 
Here, these different types of edits have a weighted cost each, and the ratio of their costs is denoted by $\alpha$. 
We then solve the purely timed alignment problem in this setting for a large class of these weighted distances (corresponding to $\alpha \in \{1\} \cup [2, \infty)$). For these distances, we provide linear time algorithms  for both distance computation and alignment on models with sequential causal processes. 
\keywords{Conformance checking \and Alignments \and Timestamps \and Time Petri nets}
\end{abstract}
%
%
%
\section{Introduction} 
    
%   \subsection{Conformance Checking and Alignments} 
    
Conformance checking %% \cite{CarmonaDSW18}
is the task of evaluating the quality of the discovered
process models with respect to event logs, i.e. to determine how well a process model
represents a system observed via the observed traces.
    
A central notion in conformance checking is the notion of an alignment \cite{Adriansyah2014AligningOA}, or, the minimal series of edits, usually insertions or deletions, that transform an observed trace into a process trace.  
    Alignments thereby help pinpoint exactly where inevitable deviations from expected behaviour occur, and the more distant the aligning word of a model is to its observed trace, the worse the model is at reflecting real system behaviour. 
    
    Assuming the event logs are a list of words over a finite alphabet (the set of possible discrete events), the problem of calculating the alignment has been extensively studied \cite{Adriansyah2014AligningOA} \cite{BCC21}, usually using either Hamming distance or Levenshtein's edit distance. 
    It is interesting to study process mining for explicitly timed systems \cite{RMAW13} \cite{CRHA20}  \cite{AS21}, that is, systems that record the timestamps that show when each event happens. 
    By considering events along with their timestamps when mining processes, we can glean information about the minimum delay between two events, the maximum duration the system takes to converge upon a state, or check deadlines, all of which are highly relevant in real-world applications \cite{Cheikhrouhou2014TheTP} \cite{Eder1999TimeCI} \cite{inbook}. 
    
    In order to study alignments for time-aware conformance checking, one requires a notion of distance over timed words, enabling us to measure how close two timed process traces are. 
Distances over timed words have been explored before, such as in \cite{ABD18}, \cite{GHJ97} and \cite{CJM14}, which capture differences between the timestamps attached to the same action label. 
In \cite{ABD18}, the distance function used can compare words of different lengths, using a Hausdorff-like pseudometric. In \cite{GHJ97}, two of the metrics defined use the stamp and delay notions we use for our edits, but restricted to linear traces. In \cite{CJM14}, a very close notion of editing first the letters and then timestamps once the letters have been aligned completely is used, but the timestamp sequences are compared by their overall deviation, so same individual deviations in position do not accrue more distance. 
Here we study and generalise the distances defined in \cite{CR22}, whose crucial feature is that timestamps are identified with the events they represent in the underlying model, rather than just the action labels. 

As for process models, there are ways to use existing process mining notation to denote time constraints. BPMN 2.0 comes equipped with \emph{timer events}. 
    For our purposes, we use time Petri nets, an extension of Petri nets equipped with the ability to express constraints on the duration of time between an action being enabled, and its actual occurrence. 
    In particular, in this paper, we restrict ourselves to time Petri nets with no branching points. 
    
    As time-aware process mining grows popular, new quality measures and conformance-checking techniques must be developed that are sensitive to temporal constraints, but so far in the study of alignments, we notice that the process model used is rarely time-aware. 
    
    Our contribution is considering a wider class of timed distances that we call weighted mixed distance functions or $d_{\alpha}$, parametrised by $\alpha$, the ratio of different costs of types of timed edit moves. 
    In doing so, we generalise the three distances defined in \cite{CR22} and in this paper solve the purely timed alignment problem for a much more general class of metrics. 
These generalised weighted mixed distances allow users to decide what types of corrections in the observed behaviour should cost more to fix, as sometimes changing an event locally is the more viable option, while in other use cases, it might be more likely that shifts in local events necessary cause consequent events to reflect the same shift. 
    
%   \subsection{Timed Alignments}
    
    \begin{example} Consider a model of composing and sending an email, where the intervals signify allowed durations between the current event and its immediate predecessor, depicted below with input and output places marked with an $i$ and an $o$: 
        
        \adjustbox{center, scale = 0.8}{
            \begin{tikzcd}
                \eye_i & {\overset{[2,3]}{\fbox{\textrm{dispatch}}}} & \bigcirc & {\overset{[5, 5]}{\fbox{\textrm{sent}}}} & \bigcirc_o \\
                {\underset{[0, 3]}{\fbox{\textrm{save draft}}}} & {\underset{[0,5]}{\fbox{\textrm{unsend}}}} & {}
                \arrow[curve={height=-12pt}, from=2-1, to=1-1]
                \arrow[curve={height=-12pt}, from=1-1, to=2-1]
                \arrow[from=1-1, to=1-2]
                \arrow[from=1-2, to=1-3]
                \arrow[from=1-3, to=1-4]
                \arrow[from=1-3, to=2-2]
                %\arrow[from=2-3, to=2-2]
                %\arrow[from=2-3, to=1-4]
                \arrow[from=1-4, to=1-5]
                \arrow[from=2-2, to=1-1]
            \end{tikzcd}
        }
        

        One process trace is $(\textrm{save draft}, 3)(\textrm{dispatch}, 5)(\textrm{sent}, 10)$, which depicts drafting the message in 3 units and then dispatching it in 2 units, and having it send successfully.
        An example of an observed trace that does not conform to the process above would be $(\textrm{save draft}, 3)(\textrm{save draft}, 5)(\textrm{dispatch}, 6)(\textrm{dispatch}, 6)(\textrm{sent}, 11)$. 
        This trace has an extraneous event as a message cannot be  dispatched twice, and the timestamps for dispatch and sending come too early. 
        One possible optimally close process trace for this observation is $(\textrm{save draft}, 3)(\textrm{save draft},4)$ $(\textrm{dispatch}, 6)(\textrm{sent}, 11)$, which deletes the extra dispatch event and prepones the second draft saving event by a unit, allowing the rest of the trace to thereby be in time. 
    \end{example}
    Here we recall the general alignment problem, posed in \cite{CR22}, which formalised the notion of aligning timed traces to timed process models.  %Time Petri nets offer both a graphical means by which to represent concurrent systems, but their semantics also allow one to express complex time constraints. 
    
    \begin{definition}[The General Alignment Problem]
        Given a process model $N$ denoted by a time Petri Net and an observed timed trace $\sigma$ we wish to find a timed word $\gamma \in \mathcal{L}(N)$ such that $d(\sigma, \gamma) = \min_{x \in \mathcal{L}(N)} d(\sigma, x)$ for some distance function $d$ on timed words. 
    \end{definition}
    
    In this paper, we restrict ourselves to the study of the alignment problem to a specific class of time Petri nets, ones that we define as sequential process models. 
    These lose the ability to express concurrent events the way time Petri nets can but retain the ability to reason about time constraints on the durations between any two consecutive events. 
    %In addition, we study a generalisation of the distances over timed words defined in \cite{CR22}, thereby considering distances where for differing costs, edits can either stay local to a given event or propagate partially forward into the event's causal future. 
    The last assumption we make is that the untimed part of the observed trace conforms to the process, i.e., the only part that requires aligning is the timestamps. 
    This allows us to focus on the timing aspect of the problem, although once this is solved we can easily adapt existing untimed alignment methods like those described in \cite{Adriansyah2014AligningOA} or \cite{BCC21} to align words with both action labels and timestamps that need editing. 
    This brings us to the following formal problem: 
    
    \begin{definition}[The Purely Timed Alignment Problem for Sequential Process Models]
        Given a sequential process model $N$ defined as a time Petri Net with no branching transitions, and an observed timed trace $(w, \sigma)$, where $w$ is the sequence of action labels and $\sigma$ is the sequence of timestamps. Suppose $w \in Untime(\mathcal{L}(N))$, i.e., the action label component of the timed word is already aligned to the model. We wish to find a valid timestamp sequence $\gamma$ such that $$d(\sigma, \gamma) = \min_{x \in \mathcal{L}(N)} d(\sigma, x)$$
    \end{definition}
    
    
    \begin{example}
    Consider the process model $N_1$: 
    
    \adjustbox{scale = 0.9}{
        \begin{tikzcd}
            &{}&& {}&&{} & \\
            \eye &\fbox{d}& \bigcirc & {\fbox{e}} & \bigcirc & \fbox{f} & \bigcirc_f\\
            \arrow[from=2-1, to=2-2]
            \arrow[ from=2-2, to=2-3]
            \arrow[from=2-3, to=2-4]
            \arrow[from=2-4, to=2-5]
            \arrow[from=2-5, to=2-6]
            \arrow[from=2-6, to=2-7]
            \arrow["{[1,3]}"{description, pos=0.2}, draw=white, from=2-2, to=1-2]
            \arrow["{[1,4]}"{description, pos=0.3}, draw=white, from=2-4, to=1-4]
            \arrow["{[0, 3]}"{description, pos=0.2}, draw=white, from=2-6, to=1-6]
        \end{tikzcd}
    }

Say we observed the trace $(u_1, \sigma_1) = (d, 4)(e, 6)(f, 6)$ that did not fit the model, and we wished to analyse how best to modify them to fit them back into the model. 
The untimed part is already aligned, and in order to make the timestamps fit we can edit the timestamp of $d$ to give $(v_1, \gamma_1) = (d, 3)(e, 6)(f, 6)$. 
    
    
    \end{example}

This paper is an extended version of our paper in the workshop proceedings of BPI 2023. 
It includes a more in depth discussion of the proof techniques used for obtaining our results. 
In particular, this includes insights on the three key properties we use to refine our edit move sequences, namely chronology, cooperation and stability. 
We believe these properties are very general, and with the appropriate extension will assist in solving the purely timed alignment problem for a wider structural class of models, and so a deeper discussion about them is valuable. 

   The rest of the paper is organised as follows.  In sections \ref{2} and \ref{3} of this paper we will focus on the problem of computing a particular case of our mixed distances when costs are equal, $d_1$ and in section \ref{4}, we return to the problem of aligning sequential process models to observed traces. 
    In subsection \ref{3} we also show that computing and aligning $d_{\alpha}$ for $\alpha \geq 2$ can be reduced to the delay only case, which was solved in \cite{CR22}. 
    We present linear time algorithms for both computing $d_{\alpha}$ and solving the alignment problem in this setting. 


    \section{Preliminaries: Edit Moves for Timed Words}
    \label{2}
    
As this paper is intended as a follow up to \cite{CR22}, we cite it as a reference for further reading about the definitions of time Petri nets or timed words in the timed alignment context. For the rest of this paper, we restrict our attention to sequential process models instead of time Petri nets, and sequences of timestamps instead of timed words. 

Inspired by Levenshtein's edit distance, we view the definition of these distances as an exercise in cost minimisation over some set of transformations between words. 
    We define \textit{moves} as functions that map one time sequence to another. 
    What are useful notions of transformation on a timed system? 

    \begin{example}
        Say we seek to align the timed trace $(4,6,6)$ with $(3, 6, 6)$. 
        It feels reasonable to say the distance between the two is 1. 
        A way to arrive at this conclusion is that if the first timestamp were just shifted to fire at 3 instead, the whole process would match. 
        This sort of local, almost typographical error can often happen in systems, and it is the simplest kind to fix.  
        
        Consider aligning the timed traces $(4, 8, 11)$ and $(3, 7, 10)$. 
        When trying to compare these firing sequences, we can view it as before, as all of the timestamps firing later than they should, and so editing one by one, getting an aligning cost of 3. 
        However, this cascading chain of errors can also be fixed if the first event is moved back to fire at 3, and all the relative relationships between it and its successors are preserved. 
        This views the second two tasks as being dependent on when the first ends, which makes sense because they are causally linked in the model. 
        Hence the switch from the timestamp series $(4, 8, 11)$ to $(3, 7, 10)$ can be viewed as only a cost 1 edit. 
        This viewpoint is of practical use as if a delay at the beginning caused the whole process to occur too late, the important thing to fix is the delay at the beginning and the rest of the process will follow suit. 
    \end{example}
    
    Based on the above example, we naturally arrived at two types of moves in \cite{CR22}.  
    We define a \textit{stamp move} as a move that translates the timing function only at a point, i.e., that edits a particular element of the timestamp series $\tau$. 
    Calculating $d_t$ between two traces can be viewed as a cost minimisation process in aligning the traces using only stamp moves. 
    
    \begin{definition}[Stamp Move]
        Given a timing function $\gamma: \{1, \dots n\} \to \mathbb{R}$, formally, we define this as: 
        
        $\forall x \in \mathbb{R}, i \leq n: stamp(\gamma, x, i) =\gamma' $ where
        $$\forall i \leq n: \gamma'(j)  = \begin{cases}
            \gamma(j) + x & j = i \\
            \gamma(j)& otherwise
        \end{cases}$$
        
    \end{definition}
    We define the cost of such a stamp move as the stamp cost per unit ($c_s$) times the magnitude of the move, i.e., $c_s \cdot |x|$. 

    The next type of move we describe is the more interesting \textit{delay move}. 
    By formulating this type of edit move, we seek to leverage the structure of the process model itself, by reflecting the causal relationships between events, yielding:
    % A stamp move is purely local, in that when a stamp edit occurs at $e$ its immediate successor events shift their relationship with $e$, thereby ensuring that the change in $e$'s timestamp does not derail the rest of the system. 
    % On the other hand, any delay move at $e$ will preserve relative relationships in the future, at the cost of shifting the timestamp of every causal descendent of $e$ by the same amount. 

    
    \begin{definition}[Delay Move]
        Given a timing function $\gamma: \{1, \dots n\} \to \mathbb{R}$, formally, we define this as: 
        
        $\forall x \in \mathbb{R}, i \leq n: delay(\gamma, x, i) =\gamma' $ where
        $$\forall i \leq n: \gamma'(j)  = \begin{cases}
            \gamma(j) + x & j \geq  i \\
            \gamma(j)& otherwise
        \end{cases}$$
        
    \end{definition}
    
    We define the cost of this delay move as the delay cost per unit ($c_d$) times the magnitude of the move, i.e., $c_d \cdot |x|$. 
The cost of a sequence of moves is the sum of the costs of the moves. Armed with these definitions, up to scaling, all the distances we can build from these depend solely upon the ratio of stamp cost to delay cost, giving: 
    \begin{definition}[Weighted mixed distance: $d_{\alpha}$]
        Given  $c_s, c_d \in \mathbb{R}^+ \cup \{\infty\}$ and assuming that we always have $c_d \cdot c_s \neq 0$, for any two timing functions $\tau_1, \tau_2$ over the same causal process $(CN, p)$ and $\alpha \in \mathbb{R}^+ \cup \{\infty\}$, we define the $\alpha$ weighted mixed distance $d_{\alpha}$ as follows 
        \small  $$d_{\alpha}(\tau_1, \tau_2) = \min \{cost(m) | \frac{c_s}{c_d} = \alpha \textrm{ and } {m \in (Stamp \cup Delay)^*}, {m(\tau_1) = \tau_2}\}$$
        
    \end{definition}
    
    From this definition, we can recover the stamp only ($d_t$), delay only ($d_{\theta}$) and evenly mixed ($d_N$) distances defined in \cite{CR22} as natural special cases, where we discount a type of move from the possible move set conceptually by setting its cost to $\infty$:
        
    $d_0(\tau_1, \tau_2) = d_t(\tau_1, \tau_2) = \min \{cost(m) | {m \in Stamp^*}, {m(\tau_1) = \tau_2}\}$

        $d_1(\tau_1, \tau_2) = \min \{cost(m) | {m \in (Stamp \cup Delay)^*}, {m(\tau_1) = \tau_2}\}$

        $d_{\infty}(\tau_1, \tau_2) = d_{\theta}(\tau_1, \tau_2) = \min \{cost(m) | {m \in Delay^*}, {m(\tau_1) = \tau_2}\}$

    Henceforward for consistency we will use $d_0, d_1,$ and $d_{\infty}$. 
    
    \begin{example}\label{stamptoy}
        Let us try to align the observed trace $(0, 3, 4)$ to the process trace $(0.5, 2.5, 3.5)$
        
        The best $d_0$ alignment has cost 1.5, with stamp moves editing each position, so $0.5$ cost for each of three edits. 
        
        The best $d_{\infty}$ alignment is cost 1.5, as it requires a 0.5 delay edit at the first place to push the 0 forward to a 0.5, and then another -1 delay at the second position to pull the rest of the trace back from $(3.5,4.5)$ back to $(2.5,3.5)$. 
        
        Now, $d_1$. 
        In sequential models, the cascading effect of a delay is not affected by the order of moves, we can align timestamps left to right. 
        Our first move must incur a minimum cost of $0.5$ as we try to align $(0,3,4)$ to our process trace at the first position.
        If any part of this initial move was a delay, we would push the second component even further away from $2.5$ than it already was, and if any part of the delay were negative, then it seems to counter the effect of the stamp part, so we should only use stamp here. 
        For move two, we again incur a minimum cost of $0.5$. 
        If the stamp part of this move were positive, it would leave the third component further off from the goal $3.5$ than if the whole move used delay, and so a pure delay move seems best. 
        This gives a cost 1 alignment, a stamp move of value 0.5 at the first position followed by a delay move of 0.5 at the second. 
        All the intuitive reasoning given above will be justified in subsequent sections, and in fact, distance 1 is the best we can do here even in the mixed case. 
    \end{example}
    
\section{Computing mixed distances}
\subsection{Notation and Setup}
    Given the nonconstructive nature of the definition of $d_{\alpha}$, it is not clear how one can efficiently calculate the distance between two fixed timed traces, as a minimal cost sequence of moves is not obvious. 
    %As before, we first study the problem over sequential causal processes, as these make analysing the effects of moves significantly simpler. 
    Before we propose an algorithm that does calculate the minimal cost for sequential timed traces, we first define a few properties that seem to characterise classes of well-formed minimal cost runs, and then prove that these properties both improve the cost and are satisfied solely by the run calculated by the algorithm we provide. 
    We start with some convenient notation for a common combination of the previously defined moves. 
    
    \begin{definition}[Mixed Move]
        We define \textit{mixed moves}, which denote doing a stamp move and a delay move at the same position in the word. 
        We define their effect as: $(s, d, e)(\gamma) = stamp(delay(\gamma, d, e), s, e)$
        Let the set of all mixed moves be $Moves$. 
    \end{definition}
    
    Given any move $m \in Moves$ we define the function $cost_{\alpha}: Moves \to \mathbb{R}^+$ that returns the cost of the move.  
    The cost of a mixed move $(s, d, e)$ is the sum of the cost of the stamp and delay moves it's made up of, $c_s\cdot|s| + c_d \cdot|d|$.
    By default by $cost((s, d, e))$ we are refering to the evenly mixed cost $cost_1((s, d, e)) = |s| + |d|$. 
We say a sequence of moves $m_1m_2\dots m_n = m \in Moves^*$ \textit{aligns} $\tau_1$ to $\tau_2$ if $m(\tau_1) = m_n(\dots (m_2(m_1(\tau_1))\dots) = \tau_2$. 
    
 Now, we look at a new way to represent timing functions. 
 %With delay edits, which constitute a new way of thinking about how a time series can be transformed, comes a new perspective with which we can view timing functions over causal processes. 
    There is of course the standard definition, $\tau: \{1, \dots n\} \to \mathbb{R}^+$ that assigns to each event a timestamp that records exactly when the event occurs. 
    Instead, thinking along the lines of delays and durations between events occurring, we define a way by which to view a timed word not in terms of its absolute timestamps, but by the delays between relevant timestamps. 
    
    \begin{definition}[Flow Function]
        Given a (not necessarily valid) time sequence,  $\tau: \{1, \dots n\} \to \mathbb{R}^+$, we first define the \textit{flow} function of $\tau$, $f_{\tau}: \{1, \dots n\} \to \mathbb{R}^+$ such that 
        $$f_{\tau}(i) = \begin{cases} 
            \tau(i) & i = 1\\
            \tau(i) - \tau(i-1) &i > 1\\
        \end{cases}
        $$
        
    \end{definition}
    
    The flow function thereby is a dual representation of timing functions, much like the relationship between graphs and line graphs. %, as if a timing function traditionally labels its transition nodes with timestamps, the flow function labels edges leading up to transitions with the duration since said transition was enabled. 
    
    \begin{example}\label{wordandflow}
        Consider the linear timed trace $w = (1, 5, 9)$. The flow function measuring its successive delays is $(1, 4, 4)$. 
    \end{example}
    
    As defined, $f_\tau$ produces exactly the time durations that the guards of each transition in the model check, i.e. the clock function values during the run.
    Hence, we see that if a word is in the language of the model, then its $f$ function maps events to values that lie within the constraint that the event's corresponding transition demands. 
    %This condition is unfortunately not sufficient though, due to urgency concerns, but it is quite close to the exact condition necessary for a word to be in the language of a time Petri net. 
    Also note that given the underlying causal process and the resulting $f_\tau$, we can reconstruct $\tau$ quite straightforwardly as $\forall i \leq n: \tau(i) = \sum_{j \leq i} f_\tau(j)$. 
    
    The following lemma characterises how moves look in flow function notation, highlighting how they are both relatively "local" in effect now:
    
    \begin{lemma}\label{FlowFuncTrans}
        Given a flow function $f_{\tau_1}$ and a mixed move $(s, d, i)$, the consequence of performing the move is the flow function defined below, where $\tau_2 = (s, d, i)\tau_1$. 
        $$f_{\tau_2}(j) = \begin{cases} 
            f_{\tau_1}(j) + s + d & j = i\\
            f_{\tau_1}(j) - s  & j = i+1\\
        \end{cases}
        $$
        
    \end{lemma}
     \begin{proof}
        This follows immediately from the definitions of a move and a flow function. 
     \end{proof}
    
    In order to compute the mixed distance between two timed traces, we wish to cull down the space of possible sequences of moves that transform the two words to each other to some move sequences that still capture all the behaviours we're interested in. 
    Hence, we introduce the following properties, which are designed intuitively to be sequences of moves that still perform the same effective transformations, and generally lower the cost of the transformation as well, by choosing moves wisely to avoid inefficiency. 
    Chronology, the first property, represents the fact that moves can commute in a sequence. This holds for all $d_\alpha$, irrespective of $\alpha$. 
    The second, co-operation, represents the fact that doing an unnecessarily expensive move in the moment will not give unintuitive future gains, and we can show this holds for all $d_{\alpha}$ where $\alpha \geq 2$. 
    The final property, stability, pinpoints exactly at what point the optimal move for the evenly mixed distance, $d_1$ is. 
    
    \subsection{Chronology}

    The first property captures the fact that we can order our moves in a sequence that goes from left to right or right to left along the positions in the sequential trace. 
    
    \begin{definition}[Chronology]
        A sequence of moves aligning $\gamma$ to $\sigma$ is said to be \textit{chronological} if for all positions $i < j \leq n = |\gamma|$, all the moves at position $i$ are performed before any move at position $j$ and exactly one mixed move (where one or both components may be zero) takes place at each position, that is $\rho \in Moves^*$ is chronological iff  $\forall i \in \{1, 2, \dots n\},  \exists s_i, d_i \in \mathbb{R}: \rho = (s_1, d_1, 1)(s_2, d_2, 2)\dots (s_n, d_n, n) $. 
        Such that $\forall i \in \{1, 2, \dots n\}: \sigma_i = \gamma_i + s_i + \sum_{j=1}^{i} d_j$. 
        
    \end{definition}
    
    \begin{definition}[Reverse Chronology]
        A sequence of moves aligning $\gamma$ to $\sigma$ is said to be \textit{reverse chronological} if for all positions $i < j \leq n = |\gamma|$, all the moves at position $i$ are performed after any move at position $j$ and exactly one mixed move (where one or both components may be zero) takes place at each position, that is $\rho \in Moves^*$ is chronological iff  $\forall i \in \{1, 2, \dots n\},  \exists s_i, d_i \in \mathbb{R}:  \rho = (s_n, d_n, n)(s_{n-1}, d_{n-1}, n-1)\dots (s_1, d_1, 1) $. 
        Such that $\forall i \in \{1, 2, \dots n\}: \sigma_i = \gamma_i + s_i + \sum_{j=1}^{i} d_j$. 
    \end{definition}
    
    Clearly, all chronological sequences have a one to one correspondence with reverse chronological sequences with exactly the same cost, the map between the two being simply reversing the sequence of moves. 
    In addition, in the flow function representation, the notions of chronology and reverse chronology are still somewhat preserved. 
    A mixed move affects both the duration before and all those after the event the move is applied at, so the order in which they happen does indeed stay chronological or reverse chronological as the case may be. 
    The only point of nuance here is that of course, a prefix of a sequence of say chronological moves no longer has the property of having completely aligned the word up to the i-th element of the flow vector while leaving the rest unchanged, as the last edit move may have perturbed the (i+1)st element. 
    
    But the upshot is, in a reverse chronological sequence of moves, the prefix of the sequence that stops at the ith event, combined with the stamp component of the mixed move at position (i-1), completely aligns the flow function for the suffix starting at position i. 
    Reverse chronological runs do not have an equivalent partial completion property in the timing function setting, but due to how local the moves end up being in the flow function setting, we obtain it here. 
    
    \textbf{Note: } For simplicity, we always assume the move played on the last position of a trace henceforth is a pure delay. Even if part of it were stamp, it would not affect the flow or timing representation any differently, so this is safe to assume. 
    
    \begin{example}
        
        Given $\gamma = (1, 1, 2, 4, 5)$ and $\sigma = (1, 2, 2.5, 4.2, 5)$, an example of a cost 5.3 non-chronological sequence of evenly mixed moves aligning them is as shown: 
         $\gamma \xmapsto{(-1, 0, 1)} (0, 1, 2, 4, 5) \xmapsto{(0, 2, 1)} (2, 3, 4, 6, 7) \xmapsto{(0, -1, 1)} (1, 2, 3, 5, 6) \xmapsto{(0.3, -0.8, 3)} (1, 2, 2.5, 4.2, 5.2) \xmapsto{(0, -0.2, 5)} \sigma$

        There is an improved cost 3.3 reverse chronological sequence of moves obtained by reordering the above sequence and combining moves at the same position: $\gamma \xmapsto{(0, -0.2, 5)} (1, 1, 2, 4, 4.8) \xmapsto{(0, 0, 4)} (1, 1, 2, 4, 4.8) \xmapsto{(0.3, -0.8, 3)}(1, 1, 1.5, 3.2, 4)\xmapsto{(0, 0, 2)} (1, 1, 1.5, 3.2, 4)\xmapsto{(-1, 1, 1)} \sigma $. 
        
    \end{example}
    
    As the above example suggests, we claim: 
    
    \begin{lemma}\label{linescommute} For all $\alpha \in \mathbb{R}^+ \cup \{\infty\}$ there is a minimal cost sequence of $\alpha$ weighted mixed moves aligning any two sequential timed traces that is chronological/reverse chronological. 
    \end{lemma}


\begin{proof}
    We assume that there is at least one mixed move at each position, adding zero moves as needed. We first observe the following: 
    $$\forall  i < j : (s, d, i)(s', d', j)(\gamma) = \gamma' = (s', d', j)(s, d, i)(\gamma)$$
    Where $$\forall 1 \leq k \leq n : \gamma'_k = \begin{cases}
    \gamma_k & k <i \\
    \gamma_k + s + d & k = i \\
    \gamma_k + d & i < k < j \\
    \gamma_k + s' + d + d' & k = j \\
    \gamma_k + d + d' & k > j
    \end{cases}$$
    That is, on sequential timed traces, mixed moves are commutative. 
    Clearly permuting the moves of a run does not change its cost, so we can assume any sequence of moves to be in increasing order of positions without loss of generality.
    
    Now we can write $m'$ as follows : 
    $$m' = (s_{11}, d_{11}, 1)\dots (s_{ij}, d_{ij}, i)\dots(s_{nk_n}, d_{nk_n}, n)$$ 
    
    Where for each $i$, $k_i \geq 1$ is the number of mixed moves at position $i$. 
    $$cost(m') = \sum_{i = i}^{n} (\sum_{j = 1}^{k_i} c_s\cdot|s_{ij}| + c_d \cdot|d_{ij}|)$$ 
    
    Define $s_i = \sum_{j = 1}^{k_i} s_{ij}$, $d_i =\sum_{j = 1}^{k_i} d_{ij}$, giving us the chronological run $$m = (s_1, d_1, 1)\dots (s_i, d_i, i)\dots (s_n, d_n, n)$$
    
    Now, by the triangle inequality we know that : 
    $$\forall i : \sum_{j = 1}^{k_i} c_s \cdot|s_{ij}| + c_d \cdot|d_{ij}| \geq c_s \cdot| \sum_{j = 1}^{k_i} s_{ij} | + c_d \cdot| \sum_{j = 1}^{k_i} d_{ij} | = cost(s_i, d_i, i) $$
    
    Thereby giving us that $$cost(m) \leq cost(m')$$ 
    
    As for the reverse chronological sequence, by commutativity we can simply reverse $m$ and obtain $m_R$, a reverse chronological sequence of moves that does the same transformation and as reversing the sequence preserves cost, $$cost(m_R) = cost(m) \leq cost(m')$$ 
    \end{proof}
    
  \subsection{Co-operation}
    
    \begin{definition}[Co-operation]
        A mixed move is said to be co-operative if its delay edit and stamp edit are in the same direction, that is, 
        $(s, d, i) \textrm{ is co-operative iff } sd \geq 0$. 
        
        A chronological (or reverse chronological) sequence of moves such as $m = m_0m_1\dots m_{n-1}$ aligning $\gamma$ to $\sigma$ is said to be \textit{co-operative} if each of its moves is co-operative. 
    \end{definition}
    
    %This property is largely unaffected by whether we're viewing the words in timing or flow representation. 
    
    \begin{example}The previous example (that was a cost 3.3 evenly mixed run, recall) can hence be further improved to become a reverse chronological co-operative cost 1.7 run as follows (on flow vector): 
        %$$\gamma \xmapsto{(0, 0, 1)} (1, 1, 2, 4, 5) \xmapsto{(1, 0, 2)} (1, 2, 2, 4, 5) \xmapsto{(0.5, 0, 3)} (1, 2, 2.5, 4, 5) \xmapsto{(0.2, 0, 4)} (1, 2, 2.5, 4.2, 5) \xmapsto{(0, 0, 5)} \sigma $$
$f_\gamma \xmapsto{(0, 0, 5)} (1, 0, 1, 2, 1) \xmapsto{(0.2, 0, 4)} (1, 0, 1, 2.2, 0.8) \xmapsto{(0.5, 0, 3)}(1, 0, 1.5, 1.7, 0.8) \xmapsto{(1, 0, 2)} (1, 1, 0.5, 1.7, 0.8) \xmapsto{(0, 0, 1)} f_\sigma $. 
        
    \end{example}
    
    Once again, one can always convert an evenly mixed non-cooperative run into an equivalent cooperative one with a smaller cost, and in fact, this holds for a large number $d_{\alpha}$ as well. 
    
    \begin{lemma}\label{linescoop} For all $\alpha \geq 2 \vee \alpha = 1$, there is a cooperative minimal cost chronological (or reverse chronological) sequence of $\alpha$ weighted mixed moves aligning any two sequential timed traces. 
    \end{lemma}
    \begin{proof}
        Given any minimal cost sequence of moves aligning $\gamma$ to $\sigma$, by Lemma \ref{linescommute} we know that there is a chronological sequence that is also minimal cost, call it $m$. 
        
        Assume $|\gamma| = |\sigma| = n$. 
        
        We proceed by induction on the first index $k$ at which $m' = (s_1, d_1, 1)\dots(s_n, d_n, n)$ uses a non co-operative move. 
        The idea is to start by replacing $m'$ with a sequence whose first move is cooperative (with perhaps minor changes in the future but no changes to the past) that is overall at most as expensive. 
        Then, by gradually keep replacing moves as one goes along, since we do not affect past moves, eventually we have a completely cooperative sequence. 
        \textbf{Extreme case} : $k = n$.
        
        This gives us $$s_n d_n < 0 \implies |s_n + d_n| < |s_n| + |d_n|$$ 
        $$\textrm{Then } m = m'|_{n-1}(0, s_n + d_n, n)$$ costs less than $m'$, aligns the words and is co-operative. 
        
        \textbf{Induction hypothesis} : Suppose we've demonstrated that there is a sequence of moves $m''$ that aligns $\gamma$ to $\sigma$, is co-operative for the first $0 \leq k-1 < n-1$ steps, is non co-operative at step $k$ and $cost(m'') \leq cost(m)$
        
        \textbf{Claim} : There is a sequence of moves $m$ that aligns $\gamma$ to $\sigma$, is co-operative for at least the first $k$ steps and $$cost(m) \leq cost(m'') \leq cost(m')$$
        
        \textbf{Proof} : Let $m'' = (s_1, d_1, 1)\dots(s_n, d_n, n)$. 
        
        We first define $m_i = m''|_{k-1}$, the initial $k-1$ length sequence of moves, and $m_f = (s_{k+2}, d_{k+2}, k+2)\dots (s_n, d_n, n)$, the final sequence of moves post the moves at positions $k$ and $k+1$. 
        Hence we have that $$m'' = m_i(s_k, d_k, k)(s_{k+1}, d_{k+1}, k+1)m_f$$
        Recall that upto scaling it is sufficient to assume that $c_s = \alpha$ and $c_d = 1$. 
        
        We know that $s_kd_k < 0$. 
        This implies one of two possibilities : 
        
        \textbf{Case 1} : $|d_k + s_k| = |d_k| - |s_k|$
        
        We define the following aligning, chronological run that is co-operative up to (and including) the $k$th step : 
        $$m = m_i( 0, d_k + s_k, k)( s_{k+1}, d_{k+1} - s_k, k+1)m_f$$
        $$cost(m'') - cost(m)  $$$$ = (\alpha \cdot |s_{k+1}| + |d_{k+1}| + \alpha \cdot |s_k| + |d_k|) - ( | d_k + s_k| + \alpha \cdot | s_{k+1}| + | d_{k+1} - s_k |) $$ 
        $$\geq \alpha \cdot (|s_{k+1} - |s_{k+1}|) + (|d_{k+1}| - |d_{k+1}|)+ (\alpha \cdot|s_k| + |d_k|  -| d_k + s_k|  - |s_k |) = (\alpha |s_k|)  \geq 0$$ 
        
        \textbf{Case 2} : $|d_k + s_k| = |s_k| - |d_k|$
        
        When $\alpha \geq 2$, we can define the same aligning, chronological run that is co-operative up to (and including) the $k$th step : 
        $$m = m_i( 0,d_k + s_k,  k)( s_{k+1}, -s_k + d_{k+1}, k+1)m_f$$
        $$cost(m'') - cost(m) $$$$= \alpha \cdot |s_{k+1}| + |d_{k+1}| + \alpha \cdot |s_k| + |d_k| - (|s_k + d_k | + \alpha \cdot| s_{k+1}| + | d_{k+1} - s_k |) $$
        $$\geq \alpha \cdot |s_k| + |d_k| - (|s_k| - |d_k| + |s_k|) = (\alpha - 2)\cdot |s_k| \geq 0$$
        
        On the other hand, when $\alpha = 1$, we define the following aligning, chronological run that is co-operative up to (and including) the $k$th step : 
        $$m = m_i( d_k + s_k, 0, k)( s_{k+1}, d_k + d_{k+1}, k+1)m_f$$
        $$cost(m'') - cost(m) $$$$= |s_{k+1}| + |d_{k+1}| + |s_k| + |d_k| - (|s_k + d_k | + | s_{k+1}| + | d_{k+1} + d_k |) $$
        $$\geq |d_{k+1}| + |s_k| + |d_k| - (|s_k + d_k | + | d_{k+1}| + |d_k |) \geq 0$$
        
        By induction we can hence conclude that there is a chronological sequence of moves $m$ that is co-operative throughout and has lower cost than $m'$, i.e., if $m'$ is minimal cost so is $m$.  
        
        Now, $m_R$ is also co-operative throughout as co-operation is a property of the individual moves constituting the sequence, and $m_R$ has exactly the same moves just in the reverse order, so there is also a reverse chronological minimal cost sequence of moves aligning $\gamma$ to $\sigma$. 
        
        \end{proof}
    For the rest of this subsection assume we are always aligning using evenly mixed distance $d_1$. 
    Similarly we have \emph{cross co-operation} which states that a stamp move $s$ at a location should ideally have opposite sign to the delay $d$ at the next position, as its contribution to the next flow component is going to be $-s$. 
    This leads us to the following definition: 
    
    \begin{definition}[Cross Co-operation]
      A reverse chronological run ${m = (0, d_n, n) \dots (s_1, d_1, 1)}$ is said to be cross co-operative if no stamp component has the same sign as the delay component of the move played at the next position, that is, $\forall i < n : s_i \cdot d_{i+1} \leq 0$. 
    \end{definition}
    
    Much like co-operation above, cross co-operation can be used to improve the cost of a run, as seen in the next lemma. 
    
    \begin{lemma}\label{crosscoop}
      For every reverse chronological co-operative evenly mixed run $m'$ aligning $\gamma$ to $\sigma$, there is a reverse chronological, co-operative and cross co-operative run $m$ aligning $\gamma$ to $\sigma$ such that $cost(m) \leq cost(m')$. 
    \end{lemma}

    \begin{proof}
        Say $m' = (0, d'_n, n) \dots (s'_1, d'_1, 1)$. 
        Let $i$ be an index at which the stamp move is not cross co-operative, i.e., $s_i \cdot d_{i+1} > 0$. 
        
        Consider a run $m''$ that is identical to $m'$, except for the moves it plays at positions $i$ and $i+1$.   
        
        Let $m''$ perform the move $(0, d_i + s_i, i)$ at position $i$ and $(s_{i+1}, d_{i+1} - s_{i}, i+1)$ for the move at position $i+1$ if $|s_i| < |d_i|$ and $(s_i - d_{i-1}, d_i + d_{i+1}, i)$ at position $i$ and $(s_{i+1}, 0, i+1)$ for the move at position $i+1$ otherwise. 
        
        Clearly, the appropriate one of either of these is cross co-operative at position $i$, while also staying co-operative and reverse chronologically aligning $\gamma$ to $\sigma$, and costs $\min(|s_i|, |d_{i-1}|)$ less that $m'$. 
        
        Doing this procedure for at most $n$ times will result in a run with every stamp move being cross co-operative and with lower cost than $m'$, that is, our required $m$. 
        \end{proof}
  \subsection{Stability}
    


    The final property we define will be largely used for reverse chronological runs on flow function representations, and it is the last improvement one can make on the cost of a sequence of moves as after enforcing this property, the sequence of moves that results is unique. 
    We prefer the flow representation for our current needs because when dealing with reverse chronological models, the flow function still retains the notion of partially aligning suffixes of the traces as we go along, provides a very useful constraint. 
    
    Moreover, we define stability in terms of flow values rather than timestamp values because the effects of mixed moves on flow vectors can be studied locally, while timestamp series get perturbed throughout their length by individual moves, and hence trying to optimise for the effect of a single mixed move involves too many variables. 
  Until further notice, consider any co-operative run to be reverse chronological, and definitions and algorithms will favour aligning flow functions rather than timing functions. 
    
  It is quite clear to see that completely aligning a flow function is equivalent to completely aligning its timing function, so this is acceptable. 
    The notion of co-operation and reverse chronology leaves only one degree of freedom so to speak in our choice of aligning moves, that is, the ratio of stamp to delay at any given position, and this optimal ratio is what stability pinpoints. % Once the stamp components of a sequence of moves has been decided, the delay components are predetermined to align the flow function from right to left. 
    Now, at any point, a delay move only affects one flow component and adds to the cost of one mixed move.
    A stamp move also only adds to the cost of one mixed move, but it has the ability to affect two flow components, hence it is natural to want to choose stamp moves in a manner that corrects both the flow components it belongs to to the best of its abilities, thereby extracting as much use as possible from as little cost, and leaving the rest up to the delay component to correct. 
    
    It is this reasoning that brings us to the property we define below as stability, which denotes this prudent choice of $s$. 
    
    \begin{definition}[Stability]
        In a reverse chronological, co-operative sequence of moves, let  ${m_i = (s, d, i)}$ be a co-operative move seeking to correct the partially aligned flow function $f_\gamma = (f_\gamma(1), \dots f_\gamma(i), f_\gamma(i+1), f_\sigma(i+2) \dots f_\sigma(n))$ (the result of the run of $m$ all the way up to and including the last stamp move) to $f_\sigma = (f_\sigma(1), f_\sigma(2), \dots f_\sigma(n))$. 
        
        Let $e_i = f_\sigma(i) - f_\gamma(i)$, and $e_{i+1} = f_\sigma(i+1) - f_\gamma(i+1)$. 
        
        We say $m_i$ is \textit{stable} if 
        $$s = \begin{cases}
            0 & e_i\cdot e_{i+1} \geq 0 \\
            e_i & \textrm{Else If }|e_i| < |e_{i+1}| \\
            -e_{i+1} & \textrm{Otherwise}
        \end{cases}$$
    \end{definition}
    
    A co-operative sequence of moves is said to be \textit{stable} if each of its moves is stable.
    Note that any stable run is also cross co-operative. 
    
        \begin{example}
        The earlier cost 1.7 example can be even further improved to be stable and have cost 1.5, thereby giving us the following:
    $f_\gamma \xmapsto{(0, -0.2, 5)} (1, 0, 1, 2, 0.8) \xmapsto{(0, -0.3, 4)} (1, 0, 1, 1.7, 0.8) \xmapsto{(0, 0, 3)} (1, 0, 1, 1.7, 0.8) \xmapsto{(0.5, 0.5, 2)} (1, 1, 0.5, 1.7, 0.8) \xmapsto{(0, 0, 1)} \sigma $
        
    \end{example}
    
    As the running example suggests, a stable chronological run (or reverse chronological run) is unique, but moreover, we claim that stability, co-operation, and chronology improve the cost of the sequence of moves.  
    
    \begin{lemma}\label{stableunique} The stable, co-operative, reverse chronological sequence of moves aligning any two sequential timed traces is unique. 
    \end{lemma}
    
  \proof{
      This is clear, as given a partially aligned flow trace the choice of $s$ to make the next move stable as defined above is deterministic, and the corresponding delay moves are determined by reverse chronology and co-operation. 
}
    
    \begin{lemma}\label{stablebest} There is a stable minimal cost sequence of moves aligning any two sequential timed traces. 
    \end{lemma}

    Before we proceed with the proof of this result, we prove a small technical lemma that will help us in the proof of Lemma \ref{stablebest}. The motivation behind the following result is essentially that the mixed distance $d_N$ respect an intuitive notion of closeness, once two words have been aligned up till the penultimate place, the one whose last duration since firing lands it closer to the desired target will naturally be the easier one to align. 
Put another way, if two words, aligning to a target word, are identical to each other in all respects but the last timestamp, then the one whose last flow value is closer to the desired target flow value, is the word that is over closer to the target under $d_N$ between the two. 

\begin{lemma}\label{cooplast}
Given two timing functions $$\gamma^x = (\gamma_1, \gamma_2, \dots \gamma_{n-1}, \gamma_{n-1} + x)$$ $$\gamma^y = (\gamma_1, \gamma_2, \dots \gamma_{n-1}, \gamma_{n-1} + y)$$ Both aligning to $\sigma = (\sigma_1, \dots, \sigma_n)$, such that $|x - (\sigma_n - \sigma_{n-1})| \leq |y - (\sigma_n - \sigma_{n-1})|$ then $$d_N(\sigma, \gamma^x) \leq d_N(\sigma, \gamma^y)$$ 
\end{lemma}

\begin{proof}
Consider any run $m^y$ aligning $\gamma^y$ to $\sigma$. 
We claim there is a run $m^x$ aligning $\gamma^x$ to $\sigma$ such that $cost(m^x) \leq cost(m^y)$. 

This proves the above lemma as then the minimal cost run aligning $\gamma^x$ to $\sigma$ would have cost at most $m^x$ corresponding to the minimal cost run aligning $\gamma^y$ to $\sigma$, so $$d_N(\sigma, \gamma^x) \leq d_N(\sigma, \gamma^y)$$

Now, we can assume without loss of generality by Lemmas \ref{linescommute} and \ref{linescoop}  that $m^y$ is chronological and co-operative, say $m^y = (s_1, d_1, 1) \dots (0, d_n, n)$. 

We construct $m^x = (s_1, d_1, 1)(s_2, d_2, 2)\dots (s'_{n-1}, d'_{n-1}, n-1)(0, d'_n, n)$ differing at most at the last two mixed moves. 
As $m^x|_{n-2} = m^y|_{n-2}$ we can also conclude that $s_{n-1} + d_{n-1} = s'_{n-1} + d'_{n-1}$ by chronology, so with co-operation we can conclude that $cost(m^y) - cost(m^x) = |d_n| - |d'_n|$. 

Hence, the problem is reduced to proving that whatever the choice of $d_n$, there is a choice of $d'_n$ that fully aligns $\gamma^x$ and has atmost equal magnitude. 

Now, $|d_n| = |\sigma_{n-1} + y - \sigma_n|$ as $m^y|_{n-1}$ fully corrected up till the penultimate place, and similarly $|d'_n| = |\sigma_{n-1} + x - \sigma_n|$, hence, by definitions of $x$ and $y$, we have the result. 

\end{proof}


Now, we can proceed with the proof of Lemma \ref{stablebest}. 
\begin{proof}

Suppose we seek to align the sequential timed trace $\gamma$ to $\sigma$, and consider a minimal cost, reverse chronological, co-operative and cross co-operative sequence of moves $m'$ that aligns $\gamma$ to $\sigma$. 
We know $m'$ exists by Lemmas \ref{linescoop} and \ref{crosscoop}. 
On the other hand, consider the unique stable, co-operative, reverse chronological sequence of moves $m$ that aligns $\gamma$ to $\sigma$, that exists by Lemma \ref{stableunique}. 

We proceed by way of induction. 
In the base case of words of length 1, the claim is clearly true as the minimal cost aligning move would just be the minimum delay move, which is also the unique stable run. 
For words of length 2, i.e., $\gamma_1\gamma_2$ aligning to $\sigma_1\sigma_2$, a reverse chronological co-operative run on the flow function is of the following form : $$m' : f_\gamma \xmapsto{(0,d_2)} (f_\gamma(1), f_\gamma(2) + d_2)$$$$ \xmapsto{(s_1,d_1)} (f_\gamma(1) + s_1 + d_1, f_\gamma(2) - s_1 + d_1) = f_\sigma$$

By reverse chronology and co-operation, we have the following equation : 
$$cost(m') = |d_2| + |s_1 + d_1| = |f_\sigma(2) - f_\gamma(2) + s_1| + |f_\sigma(1) - f_\gamma(1)|$$

Clearly the only variable in the cost is the value of $s_1$, and the stable run $m$ is the run that precisely minimises the value $|f_\sigma(2) - f_\gamma(2) + s_1|$ while maintaining co-operation, so the claim holds. 

Now, let the length of $\gamma$ be $n > 2$, and suppose the claim holds for all words of length less than $n$. 
We claim that $cost(m) \leq cost(m')$. 

The two runs must diverge at some point, as otherwise their cost is equal and the claim is obvious. There are only two main cases of interest : 

\textbf{Case 1: } The first move is the same, but the second moves are different (i.e the mixed move played at the second right-most position of the trace). 

\textbf{Case 2: }The first moves of the two sequences $m$ and $m'$ (i.e the delay played on the right-most position of the trace) are different. 



If the first move where the sequences diverge are later, say at some position $i$ in the trace, then we can restrict our attention to the $i+1$ length prefix of the word and the appropriate truncations of $m$ and $m'$, and by induction hypothesis we are done. 
So, we focus on the above cases henceforth. 


\textbf{Case 1: } Let $m = (0, d_n, n)(s_{n-1}, d_{n-1}, n-1) \dots (s_1, d_1, 1)$ and 

$m' = (0, d'_n, n)(s'_{n-1}, d'_{n-1}, n-1) \dots (s'_1, d'_1, 1)$ 

As both these runs are reverse chronological, we obtain the following equalities: $$\forall i > 1 : s_i + d_i - s_{i-1} = f_\sigma(i) - f_\gamma(i) = s'_i + d'_i + s'_{i-1}$$ 
Setting $i$ to $n$ gives $d_n - s_{n-1} = d'_n - s'_{n-1} = f_\sigma(n) - f_\gamma(n)$, which implies that  since $d_n = d'_n$, $s_{n-1} = s'_{n-1}$. 

Now consider $\gamma_s = (s_{n-1}, 0, n-1)\gamma|_{n-1}$, and 

$m_s = (0, d_{n-1}, n-1)(s_{n-2}, d_{n-2}, n-2)\dots (s_1, d_1, 1)$ and 

$m'_s = (0, d'_{n-1}, n-1)(s'_{n-2}, d'_{n-2}, n-2)\dots (s_1, d_1, 1)$. 

Clearly $m_s$ is the stable run aligning $\gamma_s$ to $\sigma|_{n-1}$ and $m'_s$ is another reverse chronological co-operative aligning sequence, and $$cost(m) = |d_n| + |s_{n-1}| + cost(m_s)$$ $$cost(m') = |d'_n| + |s'_{n-1}| + cost(m'_s)|$$ and by triangle inequality and the stability of $m$, $$|d_n| + |s_{n-1}| = |d_n - s_{n-1}| = |d'_n - s'_{n-1}| \leq |d'_n| + |s_{n-1}|$$ 
Now, by induction hypothesis on $\gamma_s, m_s, m'_s$ (they're an instance of length $n-1$ words, case 2),  we deduce $cost(m_s) \leq cost(m'_s)$, by which we can conclude that $$cost(m) \leq cost(m')$$

\textbf{Case 2:} Let $m = {(0, d_n, n)(s_{n-1}, d_{n-1}, n-1) \dots (s_1, d_1, 1)}$ and 
$m' = (0, d'_n, n)(s'_{n-1}, d'_{n-1}, n-1) \dots (s'_1, d'_1, 1)$ 

As both these runs are reverse chronological, we obtain the following equalities: $$\forall i > 1 : s_i + d_i - s_{i-1} = f_\sigma(i) - f_\gamma(i) = s'_i + d'_i + s'_{i-1}$$ 
Setting $i$ to $n$ gives $d_n - s_{n-1} = d'_n - s'_{n-1} = f_\sigma(n) - f_\gamma(n)$, which implies that having performed the first delay and stamp move of each run, one gets the following :  

$$f_\gamma \xmapsto{(0, d, n)(s, 0, n-1)} (f_\gamma(1), \dots f_\gamma(n-1) + s, f_\sigma(n))$$
$$f_\gamma \xmapsto{(0, d', n)(s', 0, n-1)} (f_\gamma(1), \dots f_\gamma(n-1) + s', f_\sigma(n))$$
Where $s = s_{n-1}$, $d = d_n$, $s' = s'_{n-1}$, and $d' = d'_n$. 

Now, let $e_n = f_\sigma(n) - f_\gamma(n)$, and $e_{n-1} = f_\sigma(n-1) - f_\gamma(n-1)$
Now, by stability of $m$, we have the following definition of $s$ : 
$$s = \begin{cases}
	0 & e_n\cdot e_{n-1} \geq 0 \\
	e_{n-1} & \textrm{Else If }|e_{n-1}| < |e_n| \\
	-e_n & \textrm{Otherwise}
\end{cases}$$

By the above definition, we see that $s = \argmin_{x \cdot (e_n - x) \geq 0} {|e_{n-1} - x|}$. 

Now by cross co-operation, we know that $s' \cdot (e_n - s') \geq 0$, so this means $$|e_{n-1} - s| \leq |e_{n-1} - s'| $$$$\implies |f_\sigma(n-1) - (f_\gamma(n-1) + s)| \leq |f_\sigma(n-1) - (f_\gamma(n-1) + s')|$$

Now, consider the words $\gamma_s = (\gamma_1, \dots \gamma_{n-2}, \gamma_{n-1} + s)$ and 
$\gamma'_s = (\gamma_1, \dots \gamma_{n-2}, \gamma_{n-1} + s')$ both aligning to $\sigma|_{n-1}$. We can apply Lemma \ref{cooplast} to these, and deduce that $$d_N(\gamma_s, \sigma|_{n-1}) \leq d_N(\gamma'_s, \sigma|_{n-1})$$

Now, by the induction hypothesis, as the rest of $m$ would have been the stable minimal cost run aligning $\gamma_s$ to $\sigma|_{n-1}$, along with stability of $m$, we have the following : 
$$cost(m) = |s| + |d| + d_N(\gamma_s, \sigma|_{n-1})$$$$ \leq |d-s| + d_N(\gamma'_s, \sigma|_{n-1}) \leq cost(m')$$

And hence we are done. 


\end{proof}

    
    
    \subsection{Computing $d_\alpha$ when $\alpha \in \{1\} \cup [2,\infty)$}\label{3}
  
The use of stability, is that it provides a unique representative for the class of minimal cost edit sequences that transform a trace to another. 
Using this fact, we prove the correctness of Algorithm \ref{MixBackwardsAlgo}
  
\begin{figure}[!t]
	\removelatexerror
	\begin{algorithm}[H]
		\caption{$d_1$ Computation Algorithm}
		\label{MixBackwardsAlgo}
		\begin{algorithmic}
			\State \textbf{Input: }  $\sigma$, $\gamma$
			\State \textbf{Output: }  $d_1(\gamma, \sigma)$
			\State $cost \gets 0$
			\State $i \gets n$
			\While {$i > 1$}
			\State $a \gets f_\sigma(i) - f_\gamma(i)$
			\State $b \gets f_\sigma(i-1) - f_\gamma(i-1)$
			
			\If{$a\cdot b \geq 0$} 
			\Comment{$(0,a,i)$} %: Stability and co-operation desire opposite delays, so $d_i$ is zero}
		\State $\gamma \gets (0, a, i) \gamma$
		
		\ElsIf{$|a| < |b|$}
		\Comment{$(-a, 0, i-1)(0, 0, i)$} %: Co-operation allows for completely correcting $i+1$ as well} 
	
	\State $\gamma \gets (-a, 0, i-1)\gamma$
	\Else
	\Comment{$(b, 0, i-1)(0, a-b, i)$}% : Stability desires a full delay, and co-operation allows no more}
\State $\gamma \gets (b, 0, i-1)(0, a-b, i) \gamma$
\EndIf

\State  $cost \gets cost + |a|$ 
\State $i \gets i-1$
\EndWhile

\State $cost \gets cost + |\gamma_1 - \sigma_1|$  \Comment $(0, |\sigma_1- \gamma_1|, 1)$

\end{algorithmic}
\end{algorithm}
\end{figure}


\begin{lemma}\label{MixAlgoCorr} Given two time sequences $\sigma = (\sigma_1, \sigma_2, \dots \sigma_n)$ and $\gamma = (\gamma_1, \gamma_2, \dots , \gamma_n)$ with sequential underlying processes,
Algorithm \ref{MixBackwardsAlgo} computes $d_1(\gamma, \sigma)$. \end{lemma}


\begin{proof}
The algorithm computes precisely the unique stable run, matching the definition of stability. 
As the distance between two words is the cost of the minimal cost sequence of edit moves that transforms one trace to the other, and any minimal cost sequence has an equivalent minimal cost stable sequence, the uniqueness of stable runs imply they are always minimal cost. 
Hence, the cost of the stable sequence of edits between $\gamma$ and $\sigma$ is the evenly mixed distance between them. \end{proof}



Now we prove a surprising fact about $d_{\alpha}$ whenever $ \alpha \geq 2$. 
Since both chronology and co-operation improve cost when computing $d_{\alpha}$, whenever the cost of the stamp move is at least twice as much as the cost of the delay move, any sequence of $\alpha$ weighted mixed moves can be converted to a purely delay move sequence with lesser or equal cost. 
Hence, assuming $c_d = 1$ and $c_s = \alpha$ (which up to scaling changes nothing), we have the following lemma: 

\begin{lemma}\label{doubleisdelay} When $\alpha \geq 2$ for any sequence of $\alpha$ weighted mixed moves aligning any two sequential timed traces, there is a sequence of delay only moves that aligns the same two traces, and has at most equal cost. 
    
\end{lemma}


\begin{proof}
	Suppose $\alpha \geq 2$, and we have two sequential timed traces $\gamma$ and $\sigma$ such that we wish to align $\gamma$ to $\sigma$ using $\alpha$ weighted mixed moves. 
	By Lemma \ref{linescoop} we know there exists a minimal cost chronological and co-operative sequence $m' = (s_1, d_1, 1) \dots (s_n, d_n, n)$ that aligns $\gamma$ to $\sigma$. 

	Consider the chronological, co-operative sequence $m = (0, s_1 + d_1)(0, s_2 + d_2 - s_1)\dots(0, s_i + d_i - s_{i-1})\dots (0, s_n + d_n - s_{n-1})$. 
	When one applies $m$ to $\gamma$, the effect on the $i$th timestamp in $\gamma$ is as follows $$m(\gamma)[i] = \gamma_i + \sum_{j = 2}^{i} (s_j + d_j - s_{j-1}) + s_1 + d_1 = \gamma_i + s_i + \sum_{j = 1}^{i} (d_j) = \sigma_i$$
This is due to the assumption that $m'$ aligned $\gamma$ to $\sigma$ correctly.

Hence, $m$ also aligns the traces correctly. Moreover, we can see that the cost of $m$ is at most that of $m'$ as shown below : 
$$cost(m) = \sum_{i=2}^{n} (|s_i + d_i - s_{i-1}|) + |s_1 + d_1| \leq \sum_{i=1}^{n} (|d_i| + 2|s_i|) \leq \sum_{i=1}^{n} (|d_i| + \alpha \cdot|s_i|) = cost(m')$$

Hence, any minimal cost aligning sequence of $\alpha$ weighted mixed moves has an equivalent delay-only sequence, so this entire class of languages behaves functionally identical to $d_{\infty}$. 
\end{proof}


As a consequence of this lemma, distance computation for this large class of distances can be done using the simple greedy algorithm described in \cite{CR22}, which runs in linear time. 

\section{Purely Timed Alignment for Sequential Models}\label{4}


Armed with the ability to compute $d_\alpha$ for a wide range of $\alpha$, we now focus on solving the purely timed alignment problem. 
To this end, we begin by defining the type of process model we will be working with. 

\begin{definition}[Sequential Process Models]
We define a sequential process model $N$ of length $n$ to be a sequence of intervals $$N = \{[a_i, b_i] | a_i \in \mathbb{R}, b_i \in \mathbb{R} \cup \{\infty\}, i \leq n\}$$
In addition, we define its language $\mathcal{L}(N)$ as follows: $\{(t_1, \dots t_n) | \forall i \leq n: t_i - t_{i-1} \in [a_i, b_i]\}$ where $t_0 = 0$. 

We depict them as follows: $\eye \longrightarrow \underset{[a_1, b_1]}{\square} \longrightarrow \bigcirc \longrightarrow \underset{[a_2,b_2]}{\square}  \dots  \underset{[a_n, b_n]}{\square} \longrightarrow \bigcirc$. 
\end{definition}

These are process models for linear strings of events. Each event $e_i$ must occur at a timestamp $t_i$ such that the number of time units between the previous event occurring and the current one belongs to the current event's time interval $[a_i, b_i]$. 

We now study an example of alignment for these processes, under three possible distances. 
    
\begin{example}\label{mixedmoveex}
Consider the below underlying sequential process model and a new observed trace $\sigma = (3, 4, 5) \not \in \mathcal{L}(N)$. 
$$\eye \longrightarrow \underset{[0,1]}{\square} \longrightarrow \bigcirc \longrightarrow \underset{[2,2]}{\square} \longrightarrow \bigcirc \longrightarrow \underset{[1,1]}{\square} \longrightarrow \bigcirc$$
The best $d_0$ alignment for the example in the diagram below is $\gamma = (1,3,4)$ with minimum cost $d_0(\sigma, \gamma) = 4$, with all events occurring as late as possible. 
The best $d_{\infty}$ alignment for the example in the diagram below is also $\gamma = (1,3,4)$, but this time with minimum cost $d_{\infty}(\sigma, \gamma) = 3$, evidenced by the move sequence $(delay(-2, 1)delay(+1, 2))$. 
And lastly, the best $d_1$ alignment for the example in the diagram below is also $\gamma = (1,3,4)$, the sequence of moves being simply one stamp and one delay move at the start, $m = stamp(-1, 1)delay(-1, 1)$, and now with minimum cost $d_1(\sigma, \gamma) = 2 < \min \{d_0(\sigma, \gamma), d_\infty(\sigma, \gamma)\}$. 
We see here that the best aligning trace can be the same, but the total cost of alignment and method of computation are very different. 
\end{example}

\begin{figure}[!t]
	\removelatexerror
	\begin{algorithm}[H]
		\caption{$d_1$ Alignment Algorithm}
		\label{MixAlignAlgo}
		\begin{algorithmic}
			\State \textbf{Input: }  $\sigma$, $N$
			\State \textbf{Output: }  $\gamma = \argmin_{x \in \mathcal{L}(N)} d_1(x, \sigma)$
			\For{$i <\in \{1, \dots, n\}$}
			\State $(a,b) \gets (Eft(t_i), Lft(t_i))$
			\State $f_{\gamma}(i)= \argmin_{x \in [a, b]} {|x - f_\sigma(i)|} $
			\State $i + + $
			\EndFor
		\end{algorithmic}
	\end{algorithm}
\end{figure}


In fact, the next theorem proves the surprising fact that while $d_1$ and $d_\infty$ are distinct distances, when aligning under them, the closest aligning trace will always coincide. 
This is to say, that the closest aligning trace for $d_1$ can be calculated by greedily editing the flow function of an observation to comply with the process model, proving the correctness of the simple and efficient method of Algorithm \ref{MixAlignAlgo}. 

\begin{theorem}\label{MixAlign} Given a sequential process model $N$ and a sequential observed trace $\sigma$, the word $\gamma \in \mathcal{L}(N)$ such that $f_{\gamma}(i)= \underset{x \in [Eft(t_i), Lft(t_i)]} {\argmin}\hspace{0.5em} {|x - f_\sigma(i)|}$ also has the property $\gamma = \underset{x \in \mathcal{L}(N)}{\argmin} \hspace{0.5em}{d_1(x, \sigma)}$. 
\end{theorem}


\begin{proof}
    
    We seek to prove that for all $\omega \in \mathcal{L}(N)$ i.e $\forall \omega : f_{\omega}(i) \in [a_i, b_i]$, $$d_N(\sigma, \gamma) \leq d_N(\sigma, \omega)$$
    By Theorem \ref{MixAlgoCorr} we can henceforth use the cost of the stable run aligning two words as their evenly mixed distance. 
    
    Let $m, m'$ be the stable runs aligning $\gamma, \omega$ respectively to $\sigma$, where: 
    $$m = (0, d_n, n)(s_{n-1}, d_{n-1}, n-1) \dots (s_1, d_1, 1)$$  
    $$m' = (0, d'_n, n)(s'_{n-1}, d'_{n-1}, n-1) \dots (s'_1, d'_1, 1)$$ 
    
    In addition, we know that $cost(m) = \sum_{i = 1}^{n} (|s_i| + |d_i|)$, and $cost(m') = \sum_{i = 1}^{n} (|s'_i| + |d'_i|)$. 
    
    We measure the cost in increments, breaking the run up at each position such that its stamp move, but not its corresponding delay move has been performed. This yields a sequence of local costs $c_i$ as defined below: $$\forall i > 0 : c_i = |s_{i-1}| + |d_i|,\,\,  c'_i = |s'_{i-1}| + |d'_i|$$ Where we let $s_0 = s'_0 = 0$ for notation's sake. 
    By doing so, we can compute the cost of perfectly aligning suffixes of the flow function. 
    
    Clearly $cost(m) = \sum_{i = 1}^{n} c_i$ and $cost(m') = \sum_{i = 1}^{n} c'_i$. 
    
    We wish to prove that $$cost(m) \leq cost(m')$$ 
    
    In fact, we shall prove something stronger. 
    We define the notion of the disadvantage of $m$, which denotes in effect the amount of distance it might have to catch up in aligning due to $m'$ having done a costlier move just prior. $$dis(i) = max(0, c'_i - c_i)$$
    
    Clearly each $c_i$, $c'_i$ and $dis(i)$ is nonnegative. 
    
    Now, we claim the following : $$\forall k > 0 : \sum_{i = k}^{n} c_i + dis(k) \leq \sum_{i=k}^{n} c'_i$$
   Then, setting $k=1$ proves that $m$ has overall lower cost than $m'$. 
    
    Let us proceed by induction on $n-k$. 
    
    For the base case, $k = n$, so the equation we seek is that $c_n + dis(n) \leq c'_n$, which is clear as $dis(n) = max(0, c'_n - c_n)$. 
    
    Now, say the claim holds for all $k > K$ for some $K < n$. 
    We seek to now prove that $$c_K + \sum_{i = K+1}^{n} c_i + dis(K) \leq c'_K + \sum_{i = K+1}^{n} c'_i$$
    
    Now, if $c_K \leq c'_K$ this is clear as $c_K + dis(K) = max(c_K, c'_K)$. 
    On the other hand, suppose $c_K > c'_K$. 
    
    By the induction hypothesis, we know that $$\sum_{i = K+1}^{n} c_i + dis(K+1) \leq \sum_{i = K+1}^{n} c'_i$$ 
    So if we show that $c_K - c'_K \leq dis(K+1)$, we are done. 
    
    Now, by reverse chronology, co-operation and cross co-operation of $m$ and $m'$, and the definition of $\gamma$, we know the following inequality holds : $$|s_{K}| + c_K = |s_K + d_K - s_{K-1}| = |f_\sigma(K) - f_\gamma(K)|$$ $$ \leq |f_\sigma(K) - f_\omega(K)| =  |s'_K + d'_K - s'_{K-1}| = |s'_{K}| + c'_K$$
    
    So, $c_K > c'_K \implies |s_K| < |s'_K|$, and moreover, $$c_K - c'_K \leq |s'_K - s_K|$$
    
    Now, by stability of $m$, we know that $s_K$ was defined as below : 
    
    $$s_K = \begin{cases}
    0 & e_K \cdot e_{K+1} \geq 0 \\
    e_K & \textrm{Else If } |e_K| < |e_{K+1}|\\
    e_{K+1} & \textrm{Otherwise}
    \end{cases}$$
    
    Where $e_K =  (f_\sigma(K) - f_\gamma(K))$ and $e_{K+1} = f_\sigma(K+1) - f_\gamma(K+1) - s_{K+1}$. 
    
    We proceed to analyse the possibilities based on the only possible stable choices of $s_K$.  
    
    Suppose firstly that $s_K = 0$. By choice of $\gamma$, $$(f_\sigma(K) - f_\gamma(K))\cdot (f_\sigma(K+1) - f_\gamma(K+1)) > 0 $$$$ \implies (f_\sigma(K) - f_\omega(K))\cdot (f_\sigma(K+1) - f_\omega(K+1)) > 0 \implies s'_K = 0$$
    But $|s'_K| >|s_K| \geq 0 $ so this first subcase is not possible. 
    
    Secondly, if $(f_\sigma(K) - f_\gamma(K)) = 0$ then  $ c_K = 0 > c'_K$, which is also impossible. 

    Thirdly, if $(f_\sigma(K+1) - f_\gamma(K+1) - s_{K+1} = 0 $ then $d_{K+1} = 0 $. 
    Now, if $d_{K+1} = 0 $ and $|s'_K| > |s_K|$, then  $$dis(K+1) = |s'_K| + |d'_{K+1}| - |s_K| \geq |s'_K|  - |s_K| \geq c_K - c'_K$$ So in this case the claim holds as needed. 
    
    Otherwise, suppose $s_K = f_\sigma(K) - f_\gamma(K)$. 
    In this case, $c_K$ is once again equal to zero, contradicting $c_K > c'_K \geq 0$. 
    
    The last possibility is that $s_K = f_\gamma(K+1) - f_\sigma(K+1) -s_{K+1}$. 
    This means that once $s_K$ is played, the $K+1$ position is perfectly aligned, so $d_{K+1} = 0$. 
    Once again, this means $$dis(K+1) = |s'_K| + |d'_{K+1}| - |s_K| \geq |s'_K|  - |s_K| \geq c_K - c'_K$$ So in this case also, the claim holds as needed. 
    
    Hence, we see that in every possible case, the claim holds for $k= K$ as well, so by induction we see that $$\sum_{i = 1}^{n} c_i + dis(K+1) \leq \sum_{i = 1}^{n} c'_i$$ From which we can conclude that $$cost(m) = \sum_{i = 1}^{n} c_i  \leq \sum_{i = 1}^{n} c'_i = cost(m')$$
    
    Hence, the result holds, and $\gamma$ is an optimal word of $\mathcal{L}(N)$ to align $\sigma$ to. 
    
    \end{proof}


\begin{theorem}
	Given a trace $\gamma$ and a sequential process model $N$, the alignment problem for weighted mixed distances $d_{\alpha}$ where $\alpha = 1 $ or $\alpha \geq 2$ can be solved in $\mathcal{O}(|N|) = \mathcal{O}(|\gamma|)$ time. 
\end{theorem}
\begin{proof}
	From Theorem \ref{MixAlgoCorr} and Lemma \ref{doubleisdelay} we know that Algorithm \ref{MixAlignAlgo} computes the correct aligning trace for this precise class of distances, and the cost of the alignment can be obtained by computing the distance to the aligning trace. 
    Since both these tasks have linear runtime complexity, overall purely timed alignment for this class of distances, to sequential process models, has been proven to be computable in $\mathcal{O}(|N|)$ time. 
	\end{proof}

\section{Implementation}
As shown in the previous sections, both distance computation and alignment tasks for a large class of mixed distances is efficiently computable, in fact the provided algorithms have been proven to run in time directly proportional to the size of the model alone. 

We have implemented the $d_1$ computation algorithm and the $d_1$ alignment algorithm for sequential process models in Python, available at \url{https://github.com/NehaRino/TimedAlignments}. 
Using varying lengths of synthetically generated sequential process models and traces to align to them, we tested both our distance and alignment algorithms for $d_1$.
These tests allow us to empirically confirm the linear relationship between model size and running time, as can be clearly seen in the run-time table below. 
These running times were obtained on an Apple M1 chip with 8 cores, running at 3.2 GHz, 8GB RAM machine.

 %\todo[inline]{Thomas: I copied this from the icpmextended paper. Can you confirm?}
\centerline{
\begin{tabular}{|c| c|}
    \hline
    Process Model Size & Running Time (seconds) \\ 
    \hline
    10 & 0.00003\\
    100 & 0.00024\\
    1000 &0.00259\\
    10000& 0.02811\\
    100000& 0.33131\\
    1000000& 3.44314\\
    \hline
\end{tabular}
}

\section{Perspectives and Conclusion}
In this paper, we introduce a generalisation of the distances over timed processes defined in \cite{CR22}, namely the weighted mixed distances. 
The principle feature of these distances (denoted by $d_\alpha$) is that depending on the cost ratio parameter $\alpha$, to align two different timed traces one will either prefer to make an edit whose effect is local to the timestamp of the current event, or its effect cascades forward to all its successors. 
We demonstrate that the problem of even computing these distances is not straightforward even when we restrict our process models to have no parallelism, i.e. the time constraint on each event only relates it to a single predecessor event at most. 
We prove that weighted mixed distances with cost ratio greater than or equal to 2 are all equivalent to the degenerate case $d_{\infty}$, i.e., when one can only edit with cascading effects, for which distance and alignment computation is solved in \cite{CR22}. 
We also solve the purely timed alignment problem for these sequential  process models for a large subset of weighted mixed distances, those where $\alpha = 1$ or $\alpha \geq 2$. . 
For the evenly mixed distance $d_1$, we solve both distance computation by providing a novel algorithm, and proving a number of properties that seem to characterise well-behaved time edit sequences. 
We also show that as for purely timed alignment under $d_1$, despite the distance function being distinct from $d_{\infty}$, the closest aligning trace for both distances coincide, hence existing algorithms can identify the aligning trace.  
As far as we know, this (along with \cite{CR22}) is the first step in conformance checking for time-aware process mining, 
and much further work can be inspired from this point. 

The first further direction is to extend the techniques used here to all process models expressible by time Petri nets, and to be able to align under any weighted mixed distance. 
As seen in this paper, even over sequential process models the problem is subtle, and while considering models that branch outwards might retain our method's efficiency, we suspect models with synchronisation between different events in particular would be tougher to tackle, due to the fact that mixed moves will cease to commute with each other chronologically. 
A perspective in this direction would be to encode the purely timed alignment problem as an ILP problem, and use generic ILP solvers to get the alignments.

It is also essential to develop practical techniques for more complex classes of models, since these will provide the foundation needed to consider extensions of time Petri nets with data or resources \cite{Availability}. 
The notions of stamp and delay we employ implicitly assume that events as represented in the process model only share resources as implied by token-flow, and more complex dependencies will render several move sequences impossible to realise in the real world. 

Secondly, further investigation into the general timed alignment problem is necessary, as our proposed approach here is rather rudimentary and can certainly be improved.
Events and timestamps are interlinked and as such should not necessarily be aligned separately, one after the other, and deciding when editing an event is less practical than its timestamp is a matter that relies heavily on the domain of application. 



Lastly, there are a number of other conformance artefacts that can be set and studied in the timed setting, such as anti-alignments \cite{CBC21}, and it would be very interesting to better develop all such conformance checking methods in a manner that accounts for timed process models.


%
% ---- Bibliography ----
%
% BibTeX users should specify bibliography style 'splncs04'.
% References will then be sorted and formatted in the correct style.
%
\bibliographystyle{splncs04}
\bibliography{Conformanceieee}





\end{document}


